
\section{$\Spec A$ 的\trZariskiTopology{}的一组基：\trDistinguishedOpen}

设 $f\in A$，定义\textbf{\trDistinguishedOpen}（distinguished open set）为
\[
D(f):=\{\, [\mathfrak{p}]\in\Spec A \;:\; f\notin\mathfrak{p} \,\}
=\{\, [\mathfrak{p}]\in\Spec A \;:\; f([\mathfrak{p}])\neq 0 \,\}.
\]
它是 $f$ 在其上不消失的\trLocus。（我私下常把这个记号写成 $D(f\neq 0)$，以提醒%
自己这一点。私底下，我也把 $f$ 的这个集合唤作「\textbf{D}oesn't-vanish set」%
（不消失集），以类比 $V(f)$ 是「\textbf{V}anishing set」（消失集）。）%
\footnote{译注：$D$ 与 $V$ 分别取自 Doesn't-vanish（不消失）与 Vanishing（消失）
的词首；这点文字游戏译不进中文，只好原样奉还。}在讨论 $\Spec A_f$ 作为 $\Spec A$
的子集时，我们已经见过这个集合。例如，我们已经观察到：\trDistinguishedOpen{}
$D(f)\subset\Spec A$ 上的\trZariskiTopology{}与 $\Spec A_f$ 上的%
\trZariskiTopology{}相一致（\cref{ex:subspace-topology-inclusions}）。

这些集合之所以重要，是因为它们构成了（\trZariski）\trTopology{}的一组特别好的%
\trBase：

\begin{exercise}[容易的练习]\label{ex:distinguished-opens-base}
证明：诸\trDistinguishedOpen{}构成（\trZariski）\trTopology{}的一组\trBase。
（提示：给定一个子集 $S\subset A$，证明 $V(S)$ 的\trComplement{}为 $\bigcup_{f\in S}D(f)$。）
\end{exercise}

下面这批练习重要而不难，能帮你找到对这个概念的感觉。

\begin{exercise}[巧妙练习]\label{ex:distinguished-cover-unit-ideal}
设 $f_i\in A$，其中 $i$ 跑遍某个\trIndexSet{} $J$。证明：%
$\bigcup_{i\in J}D(f_i)=\Spec A$，当且仅当 $(\{f_i\}_{i\in J})=A$；等价地（而且
这一形式极为有用）：存在诸 $a_i$（$i\in J$），除有限多个外均为 $0$，使得
$\sum_{i\in J}a_if_i=1$。（其中一个方向会用到「$A$ 的任何\trProperIdeal{}都含于
某个\trMaximalIdeal」这一事实。）
\end{exercise}

\begin{exercise}\label{ex:distinguished-cover-finite-subcover}
证明：若 $\Spec A$ 是诸\trDistinguishedOpen{}的一个无限并%
$\bigcup_{j\in J}D(f_j)$，那么它其实已是其中有限多个之并；换言之，存在一个有限
子集 $J'$，使得 $\Spec A=\bigcup_{j\in J'}D(f_j)$。（提示：%
\cref{ex:distinguished-cover-unit-ideal}。）
\end{exercise}

\begin{exercise}[容易的练习]\label{ex:distinguished-opens-intersection}
证明：$D(f)\cap D(g)=D(fg)$。
\end{exercise}

\begin{exercise}[重要练习；参看\cref{ex:vanish-on-vi-radical}]
\label{ex:distinguished-open-inclusion}
证明：$D(f)\subset D(g)$，当且仅当对某个 $n\ge 1$ 有 $f^n\in(g)$，当且仅当 $g$
为 $A_f$ 中的\trInvertible{}元。
\end{exercise}

\cref{ex:distinguished-open-inclusion}将被我们反复用到。你完全可以从纯代数入手
解它，但下面这条几何解读或许会有帮助。（你应当试着亲手画一张图，配合这段讨论。）
在 $\Spec A$ 内部，我们有\trClosedSubset{} $V(g)=\Spec A/(g)$，$g$ 在其上消失；
以及它的补集 $D(g)$，$g$ 在其上不消失。于是，$f$ 是这个\trClosedSubset{} $V(g)$
上的一个函数（更精确地说，是 $\Spec A/(g)$ 上的函数），而按假设，它在这个%
\trClosedSubset{}的每一点处都消失。然而，凡在环的\trSpectrum{}的每一点处都消失
的函数必是\trNilpotent{}的
（\cref{thm:nilradical-intersection-primes}）。换言之，存在某个 $n$，使得
$f^n=0$（在 $A/(g)$ 中），亦即 $f^n\equiv 0\pmod{g}$（在 $A$ 中），亦即
$f^n\in(g)$。

\begin{exercise}[容易的练习]\label{ex:distinguished-open-empty}
证明：$D(f)=\varnothing$，当且仅当 $f\in\mathfrak{N}$。
\end{exercise}

\begin{exercise}[不重要练习：环的单射映射诱导带\trDense{}像的谱映射]
\label{ex:injective-ring-map-dense-image}
证明：若 $B\hookrightarrow A$，则由此诱导的诸\trTopologicalSpace{}之间的映射
$\Spec A\to\Spec B$ 具有稠密的像。
\end{exercise}

\textbf{3.5.1.}\ \emph{关于记号的注记兼警示。}\quad 记号 $D(f)$ 将来会出现两种
变体：射影\trDistinguishedOpen{} $D_+(f)$（射影\trScheme{}上「正次数的%
\trForm（form）$f$」不消失的\trLocus，见 \S4.5.8），以及 $X_f$（\trScheme{} $X$ 上
函数 $f$ 不消失的\trLocus，见\mbox{定义~6.2.8}）。这两者与 $D(f)$ \emph{几乎}
是同一回事，但又不尽相同，所以我们出于「卫生」考虑，把这几套记号分开保留。
